\begin{theorem}[counterexample to simultaneous gluing]\hypertarget{tag-0305}{}\label{push:src:B.3.1}\label{push:B31}
\StatementRecord{push:src:B.3.1}{proved}{0305}{\SeriesRef{push:section:B.1}, \SeriesRef{push:section:B.2}, \SeriesRef{push:src:B.1.2}, \SeriesRef{push:src:B.2.1}}

The finite models of \ref{push:section:B.1} and any of the plain
$u=r^{12}$ cusp models of \ref{push:section:B.2} cannot be glued to a
proper complex space over $B\times\Delta_r$ with
the prescribed central fibres on all three base
neighbourhoods. In particular the proposed simultaneous gluing is false. This conclusion does
not rule out a different cusp model of the same
punctured $S^6$ family.
\end{theorem}


\begin{proof}
\begin{enumerate}[label=\arabic*.,leftmargin=*]
\item If a further plain cusp refinement changes the top fan,
its central fibre over an ordinary $q\ne0$ has additional
components or different multiplicities and already cannot
agree with the twelve-cycle model of \ref{push:section:B.1} on the overlap.
We may therefore assume its top fan is the standard reduced
twelve-cycle. Additional mixed-height rays lie over $q=0$
and do not change the following argument over the punctured
base.

\item Suppose such a gluing exists and put
$B^\circ=B\setminus\{0,1,\infty\}$. Normalize
its reduced central fibre over $B^\circ$. In
each of the prescribed plumbing charts this is
a proper smooth family whose fibres are the
disjoint union of twelve connected ruled
surfaces. The charts and their transition
isomorphisms therefore give a locally constant
twelve-element set of connected components.
Equivalently, the Stein factor of this
normalization is an unramified degree-twelve
cover of $B^\circ$: the algebra of fibrewise
holomorphic functions is locally $\mathcal O^{12}$,
since each proper connected ruled surface has
only constant holomorphic functions. This
description constructs the cover and requires
no conjecture about components at the missing
three points.

\item The double curves give the adjacency graph of
a twelve-cycle. Thus its transition
permutations lie in the dihedral group
$D_{12}$, the automorphism group of that graph.
We record a peripheral permutation by its
deck-identification map. Holonomy defined by
lifting a clockwise path and returning its
endpoint to the initial root slice uses the
inverse permutations throughout; the relation
and the following contradiction are unchanged.
At $0$ the monodromy is a rotation by four
steps; at $1$ it is a rotation by minus three
steps, by the explicit root deck maps in
\ref{push:src:B.1.2}. Relabelling a cycle can reverse either
orientation, but changes these only to
$\pm4$ and $\pm3$. At the cusp the
monodromy is the identity by \ref{push:src:B.2.1}.

\item The peripheral relation in
$\pi_1(B^\circ)$ is $g_1g_2g_0=1$ with
the consistently chosen clockwise loops.
In any common labelling, the product of the
two finite rotations would therefore have
to be the identity. Its possible rotation
numbers are
\[
 4+3,\quad4-3,\quad-4+3,\quad-4-3
 \equiv7,1,11,5\pmod {12}.
\]
None is zero. This contradicts the peripheral
relation. The four residues and the finite
component cycles are checked exactly by
\texttt{ab\_push\_finite\_decks.py}.
\end{enumerate}
\end{proof}
